\section{Discussion and Limitations}

The results support a narrow but useful claim.  Uniform feature dispatch is a
feature-space partitioning property, not a synonym for batch load balancing.
Recursive quantile thresholds give a direct construction for equal-mass cells,
and the held-out experiments show that this construction generalizes when the
calibration and deployment distributions match well enough.

The experiments also show why uniform usage is insufficient.  Balanced
assignment gives exact counts but low NMI.  SPR-Fixed gives clean quantile
cells but weak feature agreement in the latent-class task.  SPR-Expert is
stronger because it uses expert weights as a geometry prior.  This is not a
claim that expert weights always encode useful features: the ablations show
that noisy or random expert directions do not create meaningful cells.

Training preservation remains a separate problem.  A good initial partition can
be lost if representation drift moves many examples across thresholds.  The
SlowUpdate-SPR result is intentionally reported as a negative case: non-cheating
unlabeled recalibration can still destabilize routes if boundary updates are too
aggressive.  Future work should add explicit finite-sample tree bounds,
route-margin preservation objectives, and real vision or multimodal MoE
experiments.

\section{Conclusion}

We formulate uniform routing as a constructive feature-partition problem and
introduce Spectral Partition Routing.  The theory proves that recursive
quantile splitting gives equal-mass expert cells under continuous node
projections.  The experiments verify that quantile calibration reduces usage
imbalance on Gaussian, spherical, and covariance-flat non-spherical features.
They also show the boundary: uniform counts are not feature specialization, and
locality-only splits do not reliably recover latent classes.  The strongest
variant is expert-guided SPR, which separates the two roles cleanly: quantiles
control probability mass, while expert spectral prototypes supply feature
geometry.  Training-time preservation remains an open second-stage problem, not
a consequence of balanced initialization.
